%======================================================================
\section{Configuration lamps}\label{sec:group}
%======================================================================

This section defines the group $\Gamma$ of Theorem~A, its five generators $p,q,t,c,a$ and the $25$ relations that certify its profile, and
states Theorem~A in its precise form. The idea is visible in Figure~\ref{fig:lamps}: the lamps, copies of $S_3$, sit on configurations rather than on
points, so $N$ points carry $2^N$ lamps while every generator touches at most three points. Two of the $2^N$ lamps are compared by about $N^2$ local moves,
each filled with the Dehn function of $H_3$ at scale $N$, so the area is polylogarithmic in the number of lamps. That is what the criterion
needs. Theorem~\ref{thm:criterion} turns $M$ copies of $S_3$ whose pairwise commutators have area at most $\Lambda$ into
$\log\sigma_E(r)\ge M/16$ at $r$ of order $M\Lambda$, and it approaches its ceiling $\exp(O(r))$ only when $\Lambda$ is tiny compared with $M$.

\label{sec:toward}Put differently, the criterion needs three things at once: exponentially many sites reachable by short words, pairwise
commutators of polylogarithmic area, and an amenable ambient group. The two groups of the ladder in Section~\ref{sec:ladder} each meet two of
the three. The dyadic group $\mathcal G$ is amenable and its orbit of $0$ grows exponentially, but we cannot commute two copies at points of
opposite parity cheaply. In the subgroup $\mathcal L$ of Thompson's group $V$ everything works except amenability, which is the question
whether $F$ is amenable. The configuration lamps meet all three by changing what a site is. The sites are the finitely supported binary
configurations of the points of Houghton's three rays, not the points themselves. Affine maps that touch a bounded number of coordinates at a
time carry the copies between sites. And the ambient group is locally finite-by-$\Z^2$.

\paragraph{Houghton's group.} Let $Y=\{1,2,3\}\times\{1,2,\dots\}$; the point $(i,j)$ lies on ray $i$ at height $j$. Houghton's group
$H=H_3$ consists of the permutations $g$ of $Y$ that are eventually translations on each ray: there are integers $m_1(g),m_2(g),m_3(g)$ with
$g(i,j)=(i,j+m_i(g))$ for all large $j$. The map $g\mapsto(-m_2(g),-m_3(g))$ is a surjective homomorphism $H\to\Z^2$ whose kernel is the group
$\FSym(Y)$ of finitary permutations~\cite{Houghton1978}. Let $p$ be the permutation with $p(1,j)=(1,j+1)$, $p(2,1)=(1,1)$ and
$p(2,j)=(2,j-1)$ for $j\ge2$, fixing ray~$3$, and let $q$ be defined in the same way with rays $2$ and $3$ exchanged. They generate $H$
(Appendix~\ref{sec:houghton}), and their images in $\Z^2$ are $(1,0)$ and $(0,1)$.

\paragraph{The group.} Let $V=\F_2^{(Y)}$ be the space of finitely supported functions $Y\to\F_2$, the \emph{configurations},\footnote{This $V$ is
unrelated to Thompson's group $V$, which appears only in Sections~\ref{sec:3V} and~\ref{sec:thompson}.} with basis $(e_y)_{y\in Y}$. The group $H$ acts on $V$ by permutation matrices, $g(e_y)=e_{g(y)}$, and we identify $H$ with its image in $\GL(V)$. Let
$\GL_{\rm fin}(Y)$ be the group of linear automorphisms of $V$ that fix all but finitely many basis vectors. It is the directed union of the finite
groups $\GL(\F_2^J)$ over finite $J\subset Y$, each extended by the identity on the other basis vectors. It is normalized by $H$, and
$H\cap\GL_{\rm fin}(Y)=\FSym(Y)$. Let $L=\GL_{\rm fin}(Y)\,H\le\GL(V)$. Write $\tau_u$ for the translation $v\mapsto v+u$ of $V$, and let
$A=V\rtimes L$ be the group of affine maps $v\mapsto g(v)+u$ with $g\in L$ and $u\in V$. Our group is
\[
\Gamma=\Bigl(\bigoplus_{v\in V}S_3\Bigr)\rtimes A,
\]
where $A$ permutes the summands through its action on $V$. The copy of $S_3$ at $v$ is the \emph{lamp} at $v$; writing $\lambda_v$ for the copy
at $v$ of $\lambda\in S_3$, we have $g\lambda_vg^{-1}=\lambda_{g(v)}$ for $g\in A$. Thus $\Gamma$ is a lamplighter group. The lamps are copies of
$S_3$, one at each finite configuration of lit points of the three rays. The lighter is the affine group $A$, which moves configurations by
permuting points, toggling a point, or adding one point into another. Equivalently, $\Gamma$ is the group of permutations of $V\times S_3$
generated by the maps $(v,s)\mapsto(g(v),s)$ for $g\in A$ and $(v,s)\mapsto(v,\lambda(v)s)$ for finitely supported $\lambda:V\to S_3$. This
action is faithful. Groups obtained from a permutation action by adjoining the finitary linear or symmetric groups of the permuted set, and lamps
over it, appear as enrichments~\cite{Bradford2022,AlekseevBradford2026} and as halo products~\cite{GenevoisTessera2024}.

For distinct points $y,y'\in Y$ let $E_{y\leftarrow y'}$ be the transvection that sends $e_{y'}$ to $e_{y'}+e_y$ and fixes the other basis vectors.
For $g\in H$ and a linear map $M$ of $V$,
\begin{equation}\label{eq:affine-identities}
g\tau_{e_y}g^{-1}=\tau_{e_{g(y)}},\qquad gE_{y\leftarrow y'}g^{-1}=E_{g(y)\leftarrow g(y')},\qquad M\tau_uM^{-1}=\tau_{M(u)} .
\end{equation}
Linear maps fix $0$, so they commute with the lamp at $0$. On the first ray we abbreviate: for $x\ge1$ we write $x$ for the point $(1,x)$,
$e_x$ for $e_{(1,x)}$ and $\tau_x$ for $\tau_{e_x}$.

Let $a$ and $b$ be the transpositions $(1\,2)$ and $(2\,3)$ in the lamp at $0$, so that $ab$ has order three, and put
\[
t=\tau_1,\qquad d=E_{2\leftarrow1},\qquad c=d\,ab ;
\]
thus $t$ adds the point $1$ to a configuration, and $d$ adds the point $2$ to every configuration that contains the point $1$.

\begin{lemma}\label{lem:structure}
\begin{enumerate}[label=(\alph*),itemsep=1pt]
\item The element $c$ has order six, $c^3=d$ and $c^4=ab$; hence $b=ac^4$.
\item $\Gamma$ is generated by the five elements $p,q,t,c,a$.
\item Let $\pi:\Gamma\to\Z^2$ be the composite of $\Gamma\to A\to L\to L/\GL_{\rm fin}(Y)\cong H/\FSym(Y)\cong\Z^2$. Its kernel
$\bigl(\bigoplus_VS_3\bigr)\rtimes\bigl(V\rtimes\GL_{\rm fin}(Y)\bigr)$ is locally finite. Hence $\Gamma$ is locally finite-by-$\Z^2$, and in
particular elementary amenable.
\end{enumerate}
\end{lemma}

\begin{proof}
(a) The linear map $d$ fixes $0$, so it commutes with $ab$; since $d$ has order two and $ab$ order three, $c$ has order six, $c^3=d^3(ab)^3=d$ and
$c^4=d^4(ab)^4=ab$.

(b) Let $\Gamma'$ be the subgroup generated by the five elements. It contains $H$, because $H=\langle p,q\rangle$. Since $\FSym(Y)$ acts
$2$-transitively on $Y$, \eqref{eq:affine-identities} shows that $\Gamma'$ contains every $\tau_{e_y}$, hence $V$, and, since $d=c^3$, every
transvection $E_{y\leftarrow y'}$. Transvections generate every $\GL(\F_2^J)$, so $\GL_{\rm fin}(Y)\le\Gamma'$, and $A\le\Gamma'$. Finally
$b=ac^4\in\Gamma'$, and the lamp at $u\in V$ is $\tau_u\langle a,b\rangle\tau_u^{-1}$.

(c) Since $H\cap\GL_{\rm fin}(Y)=\FSym(Y)$, we have $L/\GL_{\rm fin}(Y)\cong H/\FSym(Y)$, and the kernel of $\pi$ is as stated. It is an
extension of $\bigoplus_VS_3$ by $V\rtimes\GL_{\rm fin}(Y)$, the directed union of the finite groups $\F_2^J\rtimes\GL(\F_2^J)$; both are locally
finite. An extension of a locally finite group $K$ by a locally finite group is locally finite: a finitely generated subgroup $\Lambda$ has finite
image modulo $K$, so $\Lambda\cap K$ has finite index in $\Lambda$, is finitely generated, and is finite. Locally finite groups are directed
unions of finite groups, and elementary amenable groups are closed under directed unions and extensions~\cite{Chou1980}.
\end{proof}

\subsection{A certificate of 25 relations}\label{sec:certificate}

All fillings in $\Gamma$ are built from one finite set $R$ of words in the five generators of Lemma~\ref{lem:structure}, each representing $1$
in $\Gamma$. The set $R$ does not present $\Gamma$ (Section~\ref{sec:fp}). Areas are always taken with respect to $R$.

\paragraph{Abbreviations.} Put
\[
s=[p,q],\qquad B=psp^{-1},\qquad t_m=p^{m-1}\,t\,p^{1-m}\quad(m=1,2,3),\qquad d=c^3,\qquad b=ac^4 .
\]
By Lemma~\ref{lem:transpositions}, $s$ and $B$ are the transpositions $(1\ 2)$ and $(2\ 3)$, and $t_m$ represents $\tau_m$, since
$p^{m-1}(1,1)=(1,m)$. By Lemma~\ref{lem:stabilizers}, with $\zeta_1=s$ and $\zeta_2=Bs$, the pairs
\[
K_1=(sp,\ sq),\qquad K_2=(Bsp,\ Bsq)
\]
generate the pointwise stabilizers $H^{(1)}$ of the point $1$ and $H^{(2)}$ of the points $1,2$. The elements $t$ and $d$ have the
\emph{prototypes} $P_t=\{1\}$ and $P_d=\{1,2\}$, the points on which they depend: by~\eqref{eq:affine-identities}, an element of $H$ that fixes
$P_f$ pointwise commutes with $f$.

\begin{table}[ht]
\centering\small
\begin{tabular}{@{}l p{0.36\textwidth} p{0.29\textwidth} r r@{}}
\toprule
 & relations & says & number & length\tabularnewline
\midrule
\textsf{H} & \raggedright $s^2$,\ $(sps\,p^{-1})^3$,\ $[s,p^2sp^{-2}]$,\ $psp^{-1}qs^{-1}q^{-1}$ & \raggedright Houghton's presentation~\eqref{eq:PH} & 4 & 64\tabularnewline
\textsf{Q} & \raggedright $t^2$ & \raggedright a translation is an involution & 1 & 2\tabularnewline
\textsf{S} & \raggedright $[t,k]$ for $k\in K_1$;\ \ $[d,k]$ for $k\in K_2$ & \raggedright the stabilizer of a prototype commutes with it & 4 & 56\tabularnewline
\textsf{Z} & \raggedright $[g,x]$ for $g\in\{p,q\}$, $x\in\{a,b\}$;\ \ $[d,a]$ & \raggedright linear maps fix the lamp at $0$ & 5 & 40\tabularnewline
\textsf{C} & \raggedright $[t_1,t_2]$ & \raggedright translations commute & 1 & 8\tabularnewline
\textsf{A} & \raggedright $d\,t_1d^{-1}(t_1t_2)^{-1}$,\ $[d,t_2]$,\ $[d,t_3]$ & \raggedright how $d$ moves the translations at $1,2,3$ & 3 & 39\tabularnewline
\textsf{L} & \raggedright $a^2$,\ $b^2$,\ $(ab)^3$ & \raggedright the lamp at $0$ is $S_3$ & 3 & 30\tabularnewline
\textsf{P} & \raggedright $[x,tyt^{-1}]$ for $x,y\in\{a,b\}$ & \raggedright the lamps at $0$ and at $e_1$ commute & 4 & 64\tabularnewline
\midrule
 & total & & 25 & 303\tabularnewline
\bottomrule
\end{tabular}
\caption{The certificate $R$: eight families of relations in the five generators $p,q,t,c,a$, written with the abbreviations $s$, $B$, $t_m$,
$d$ and $b$. Each relation is taken to be the freely reduced word obtained by expanding the abbreviations; lengths are those of these words, and
the longest has length $24$.}
\label{tab:certificate}
\end{table}

\begin{lemma}\label{lem:validity}
Every word of Table~\ref{tab:certificate} represents~$1$ in $\Gamma$.
\end{lemma}

\begin{proof}
\textsf{H}: the words of~\eqref{eq:PH} hold in $H$.
\textsf{Q}: $V$ has exponent two.
\textsf{S}: the elements of $H^{(1)}$ fix the point $1$, so they commute with $t=\tau_1$; those of $H^{(2)}$ fix the points $1$ and $2$, so they
commute with $d=E_{2\leftarrow1}$ (Lemma~\ref{lem:structure}(a) and~\eqref{eq:affine-identities}).
\textsf{Z}: $p$, $q$ and $d$ are linear, so they commute with $a$ and with $b=ac^4=a\cdot ab$, which lie in the lamp at~$0$.
\textsf{C}: translations commute.
\textsf{A}: $dt_md^{-1}=\tau_{d(e_m)}$, and $d(e_1)=e_1+e_2$, while $d$ fixes $e_2$ and $e_3$.
\textsf{L}: these are relations of $S_3$.
\textsf{P}: $tyt^{-1}$ lies in the lamp at $e_1$, which commutes with the lamp at~$0$.
\end{proof}

A program in the accompanying files checks the $25$ words independently of Lemma~\ref{lem:validity}, in the faithful action of $\Gamma$ on
$V\times S_3$; Appendix~\ref{app:provenance} says what it checks.

\subsection{Theorem A}\label{sec:theorem-A}

\begin{theorem}[Theorem A]\label{thm:main-group}
Let $E\subset\Gamma$ be the chunk consisting of $1$, the five generators $p,q,t,c,a$ and their inverses, the elements represented by the prefixes
of the $25$ freely reduced words of Table~\ref{tab:certificate}, and the inverses of these elements. Let $p_H=\log107+4$ and
\[
\kappa=2+p_H=6+\log107<12.75 .
\]
There are absolute constants $c,C>0$ such that for every $r\ge2$:
\begin{enumerate}[label=(\alph*),itemsep=1pt]
\item $\log\sigma_E(r)\ge c\,r/(\log r)^\kappa$;
\item $\log\sigma_E(r)\le C\,r\log r$; more generally, every chunk $E'$ of $\Gamma$ satisfies $\log\sigma_{E'}(r)\le C_{E'}\,r\log r$.
\end{enumerate}
In particular $\sigma_E(r)=\exp(r^{1+o(1)})$. Moreover:
\begin{enumerate}[label=(\alph*),itemsep=1pt,resume]
\item For every $0<\Delta\le1$ there are $c_\Delta,e_\Delta>0$ such that the following holds for $0<e\le e_\Delta$. Let $\phi$ assign unitaries in
$U(D)$ to the five generators so that every word of Table~\ref{tab:certificate} is mapped within $e$ of $I$ in the normalized Hilbert--Schmidt
norm, and put $\phi(b)=\phi(a)\phi(c)^4$. If $\|\phi(a)\phi(b)-I\|_2\ge\Delta$, then
\[
\log D\ \ge\ c_\Delta\,\frac{e^{-2}}{(\log(1/e))^{2\kappa}},\qquad 2\kappa<25.49 .
\]
\item $\Gamma$ is finitely presented and elementary amenable.
\end{enumerate}
\end{theorem}

The exponent $\kappa=2+p_H$ has two parts. Two of the $2^N$ lamps are brought together by about $N^2$ local moves
(Theorem~\ref{thm:pair-area}), each priced by the Dehn function of $H_3$ at scale $N$, which is at most $CN^{p_H}$
(Proposition~\ref{prop:dehn}(b)). The $107$ is the factor by which the doubling endomorphism of $H_3$ in~\cite[Section~3]{DouglasMghirbiA}
multiplies a reweighted area at each step, so that far commutators at distance $M$ have area $O(M^{\log107})$~\cite[Theorem~3.2]{DouglasMghirbiA}.
The $4$ comes from reducing the Dehn function to far commutators~\cite[Theorem~3.3]{DouglasMghirbiA}. The $107$ is the only constant in
$\kappa$ that is not structural, and a quadratic Dehn function for $H_3$ would give $\kappa=4$.

The proof occupies Sections~\ref{sec:fillings}--\ref{sec:more}.
\begin{enumerate}[label=(\roman*),itemsep=1pt]
\item Section~\ref{sec:fillings}: every local identity among transported translations, transvections and lamps is filled at a cost $L(n)$
governed by $\delta_H$ at scale $n$.
\item Section~\ref{sec:pair-areas}: two of the $2^N$ lamps on the first $N$ points commute at the cost of about $N^2$ local fillings
(Theorem~\ref{thm:pair-area}).
\item Section~\ref{sec:profile}: feeding these pair areas into Theorem~\ref{thm:criterion} gives the lower bounds (a) and (c).
\item Section~\ref{sec:profile}: models that hash configurations by random affine maps give the upper bound (b).
\item Section~\ref{sec:more}: finite presentation, part (d), by Cornulier's criterion for permutational wreath products.
\end{enumerate}
